\section{Introduction}

A proposed matching can guide an online allocation rule, but the proposal may
be wrong about edges that have not yet appeared. Robustness measures performance
for every such proposal; consistency measures performance when the proposal is
an optimum matching. We determine their exact tradeoff for unweighted bipartite
graphs of maximum degree two on both sides. This restriction separates paths
from cycles and makes the cost of using advice explicit: perfect consistency
permits robustness $1/2$ when paths are allowed, and $3/4$ when every component
is a cycle.

The structural reason for the stronger cycle guarantee is local. Once an
offline vertex has appeared twice, it cannot be used by a future arrival.
Prioritizing that vertex over a fresh one protects future options. We show that
advice can break the remaining ties without sacrificing optimality under
\emph{any} maximum advised matching, even when it is partial. On cycles, every
early choice creates one occupied and one free boundary. A later failure
consumes two occupied boundaries, so arbitrary advice still leaves a
deterministic $3/4$ matching. This advice-preserving property is the main step
in the tradeoff characterization.

\paragraph{Model and guarantees.}
Let $U,V$ be labelled sets of size $n$, with $U$ offline and $V$ online.
The realized graph $G=(U,V,E)$ is simple and unweighted. Each $v\in V$ arrives
once in an arbitrary order $\pi$, revealing its whole neighborhood. An online
algorithm irrevocably matches $v$ to an unused neighbor, or leaves it unmatched.
Its output is an integral matching. The algorithm knows $U,n$, and the online
identifiers used in the advice, but does not know future neighborhoods, their
order, or the component decomposition.

The complete advice is a table $A:V\to U\cup\{\bot\}$ with distinct
non-$\bot$ values. It specifies a partial matching in the complete bipartite
graph on the given labels. An advised pair need not be an edge of $G$.
There are no annotations or additional graph information in the advice.
The graph, order, and table are fixed independently of the algorithm's private
randomness: the adversary is oblivious. Write $\OPT(G)$ for maximum matching
cardinality. An algorithm is $c$-consistent if, for every $G,\pi$ and
\emph{every} table $A$ that is a maximum matching of $G$,
\[
  \E|\ALG(G,\pi,A)|\ge c\OPT(G).
\]
It is $r$-robust if the same inequality with $r$ holds for every allowed table.
All guarantees are uniform over finite $n$ and have no additive loss; inputs
with $\OPT(G)=0$ impose no condition.

Let $\D$ contain the graphs of maximum degree at most two on both sides,
including isolated vertices. Let $\C\subset\D$ contain the 2-regular graphs,
which are disjoint unions of simple even cycles, including four-cycles.
For a class $\mathcal X$, let $r_{\mathcal X}(c)$ be the largest robustness
attainable by a $c$-consistent randomized online algorithm on that class.

\begin{theorem}\label{thm:frontiers}
For every $c\in[0,1]$,
\[
 r_{\D}(c)=\min\{3/4,3/2-c\},\qquad
 r_{\C}(c)=\min\{7/8,7/4-c\}.
\]
Both values are attained in the model above.
\end{theorem}

\paragraph{Closest results and the contribution.}
F\`abrega and Mu\~noz~\cite{FM18}, Propositions~1--2, already establish the
advice-free $7/8$ optimum on cycles, the eight-cycle obstruction used below,
and the edge-orientation argument underlying the cycle analysis. Their ratio
uses the reciprocal convention. Feng, Li, Wu, and Zhou~\cite{FLWZ25},
Algorithm~2 and its subsequent discussion, identify the relevant priority rule
as a degree-two semi-OCS rule. We credit those baselines and the orientation
method. Our proof includes them to make the path extension and the effect of
advice self-contained.

Choo, Gouleakis, Ling, and Bhattacharyya~\cite{CGLB24}, Section~3, already give
the Z-path obstruction and the linear bound $c+r\le3/2$ in the broader
adversarial-arrival setting. Jin and Ma~\cite{JM24}, Proposition~3.1, also use
this obstruction and obtain the same numerical unweighted frontier in a
two-stage model. Their algorithm sees arrivals in batches and their
consistency benchmark uses the value of following first-stage advice and
completing later. Here every online choice is irrevocable before the next
neighborhood appears, and consistency concerns every optimum full matching
table. The contribution on $\D$ is the sequential achievability argument;
on $\C$ it is the advice-preserving endpoint and the resulting exact frontier.
Convex mixing itself is standard.

Two related distinctions matter. Bhangale, Chakraborty, and
Harsha~\cite{BCH25} bound only the \emph{online} degree in their degree-two
integral theorem; their separate $3/4$ theorem allows fractional actions.
Lindermayr, Megow, and Paul~\cite{LMP26} study multistage allocation and also
give a one-by-one tradeoff for vertex-weighted \emph{fractional} matching.
An integral prediction in that model does not make its actions integral.
These statements therefore address different algorithmic classes from
Theorem~\ref{thm:frontiers}.

\section{Priority rules and the advice-preserving endpoint}
Maintain $h(u)$, the number of revealed incident edges at $u\in U$.
On arrival of $v$, increment $h(u)$ for every $u\in N(v)$.
Call a neighbor \emph{fresh} if $h(u)=1$, and \emph{exhausted} if $h(u)=2$.
An exhausted neighbor has no future incident edge. A neighbor is available
when it is unmatched. Fix a label order for deterministic ties.

Both rules below first take an available exhausted neighbor, if one exists,
and leave $v$ unmatched only when no neighbor is available.
\begin{itemize}
\item $\Fair$ ignores the table. Without an available exhausted neighbor, it
      takes the sole available neighbor, or chooses fairly between two fresh
      neighbors, independently at each such choice.
\item $\Trust$ uses the same priority. Without an available exhausted neighbor,
      it takes $A(v)$ if that is an actual available neighbor; otherwise it
      takes the first available neighbor in the fixed order.
\end{itemize}
These rules use only the revealed prefix and the matching table. In
particular, an absent advised edge is simply ineligible.

\begin{lemma}\label{lem:advice}
For every maximum matching table $A$ of a graph in $\D$, the output of
$\Trust$ has cardinality $\OPT(G)$.
\end{lemma}
\begin{proof}
Let $M$ be the output and suppose $|M|<|A|$. The symmetric difference of
$A$ and $M$ contains an $M$-augmenting path
\[
 v_0,u_0,v_1,u_1,\ldots,v_k,u_k,
\]
where $v_i u_i\in A$, $v_{i+1}u_i\in M$, and $v_0,u_k$ are unmatched in
$M$. If $k=0$, greediness gives an immediate contradiction. Otherwise $u_0$
was already occupied when $v_0$ arrived, since $v_0$ remained unmatched.
Thus its partner $v_1$ arrived earlier than $v_0$.

Inductively suppose $v_i$ arrived before $v_{i-1}$, for $i\ge1$.
The two neighbors of $u_{i-1}$ are precisely $v_i$ and $v_{i-1}$, by the
offline degree bound, so $u_{i-1}$ was fresh when $v_i$ selected it.
Consequently no exhausted neighbor was available then. If its advised
neighbor $u_i$ had been free, $\Trust$ would have preferred $u_i$ to
$u_{i-1}$: either by exhausted priority or by the advice tie-break.
Thus $u_i$ was already occupied. For $i<k$ its partner $v_{i+1}$ arrived
earlier still. At $i=k$ this is impossible, since $u_k$ is free even at
termination. The resulting strictly decreasing arrival times contradict the
existence of the augmenting path.
\end{proof}

The argument applies to every maximum table, including partial ones. The
offline degree bound, rather than perfection or uniqueness of the advice,
is what makes the selected neighbor fresh along the augmenting path.

\section{Boundary accounting on paths and cycles}
We analyze final components only for the proof; the algorithm never needs
to identify them. Their arrivals can be arbitrarily interleaved.

Within a nontrivial component, classify a degree-two arrival as a
\emph{root}, \emph{extension}, or \emph{merge} according as it has zero,
one, or two previously seen neighbors. Let their counts be $a,b,z$.
A degree-one online vertex is an \emph{initial leaf} if its neighbor is
fresh, and a \emph{terminal leaf} otherwise; let their counts be $s,t$.
These events are determined by the fixed graph and order.

An offline vertex seen once is a boundary of the revealed component. Mark
it with bit $1$ if occupied and bit $0$ if free. A root creates two
complementary boundary bits. Under $\Fair$ they are assigned by that root's
independent fair coin. An initial leaf creates one deterministic bit $1$.
At an extension the old neighbor becomes exhausted. If free, it is chosen
and the new boundary stays free; if occupied, the new neighbor is chosen
and the new boundary is occupied. Thus the extension always succeeds and
transports the old bit unchanged.

A merge consumes its two boundaries and fails exactly when both bits are
$1$. Its choice between free exhausted neighbors cannot affect a future
boundary. A terminal leaf consumes one bit and fails exactly when it is
$1$. Every surviving boundary is therefore an original root or initial-leaf
bit transported without copying. Its source is fixed by the graph and
order, independently of all matching outcomes.

Contract extensions. For a final path, the resulting graph of roots and
initial leaves linked by merges is a path, with terminal leaves or final
offline endpoints at any remaining ports. For a final cycle it is a cycle
of roots. In particular, a merge between distinct root objects under
$\Fair$ fails with probability $1/4$, and a merge between an initial leaf
and a root fails with probability $1/2$. These are unconditional
probabilities on a fixed input. No independence conditional on earlier
successes is needed.

\begin{lemma}\label{lem:expectations}
Let a nontrivial path have $m$ online vertices and $\ell\in\{0,1,2\}$
online endpoints. If $a\ge1$, then
\begin{equation}\label{eq:path}
 \E|\Fair(\text{path})|=m-\frac{a-1}{4}-\frac{\ell}{2}.
\end{equation}
If $a=0$, $\Fair$ is optimal on the path. On a cycle with $m$ online
vertices,
\begin{equation}\label{eq:cycle}
 \E|\Fair(\text{cycle})|=
 \begin{cases}m,&a=1,\\m-a/4,&a\ge2.\end{cases}
\end{equation}
\end{lemma}
\begin{proof}
For a path with $q$ offline vertices, count first appearances and all
arrivals to obtain
\begin{equation}\label{eq:counts}
 q=m+1-\ell=2a+b+s,\quad s+t=\ell,\quad
 m=a+b+z+s+t,\quad z=a+s-1.
\end{equation}
Suppose $a\ge1$. In the contracted path, an initial-leaf object has only
one port. Joining two such objects would complete an entire component
without a root, so each of the $s$ initial leaves is joined to a root.
The remaining $z-s=a-1$ merge edges join distinct roots. Likewise a
terminal leaf consumes a root bit: connecting it directly, after
extensions, to an initial leaf would complete a separate rootless path.
The expected failures are therefore
$(a-1)/4+s/2+t/2$, giving~\eqref{eq:path}.

If $a=0$, the first arrival is an initial leaf and every boundary is
occupied. Roots and extensions cannot fail; all $z$ merges and $t$
terminal leaves fail. There are $z+t=\ell-1$ failures, giving
$m-\ell+1=\min\{m,q\}$ matches, the path optimum. Isolated vertices
contribute zero and are handled separately.

For a cycle, $s=t=0$, $m=q=2a+b$, and $z=a$.
If $a=1$, the closing merge sees the complementary bits of its sole root
and succeeds. If $a\ge2$, every edge of the contracted root cycle joins
distinct roots: a self-join would already close a separate final cycle.
Every merge then fails with probability $1/4$. With two roots the
contracted cycle has two parallel edges, each using different ports.
Linearity of expectation still gives~\eqref{eq:cycle}, without requiring
independence between merge failures.
\end{proof}

\begin{corollary}\label{cor:fair}
The rule $\Fair$ guarantees $3\OPT/4$ on $\D$ and $7\OPT/8$ on $\C$.
\end{corollary}
\begin{proof}
Rootless paths are optimal. For a path with $a\ge1$ and $\ell=0$,
$\OPT=m=2a+b-1$; its deficit is $(a-1)/4\le m/8$ because
$2(a-1)\le m$. For $\ell=2$, $\OPT=q=m-1=2a+b+s$ and its deficit
relative to $q$ is $(a-1)/4\le q/8$. Thus both unbalanced path types
have ratio at least $7/8$.

For $\ell=1$, $\OPT=m=2a+b+s$. Its deficit is $(a+1)/4$, which is
at most $m/4$ since $m\ge2a\ge a+1$. On a cycle, the deficit is at most
$a/4\le m/8$, since $m=2a+b$. Sum these bounds over final components.
Interleaving does not change their induced orders or their contributions.
\end{proof}

\begin{lemma}\label{lem:robust}
For every allowed advice table, $\Trust$ produces at least $\OPT/2$ matches
on $\D$ and at least $3\OPT/4$ matches on $\C$, deterministically.
\end{lemma}
\begin{proof}
The output is maximal: an edge between two finally unmatched vertices
would have been available when its online endpoint arrived. Every optimum
edge touches an output edge, and each output edge can account for at most
two optimum edges. This proves the $1/2$ bound.

On a cycle the boundary accounting holds deterministically for $\Trust$.
Each of its $a$ roots creates exactly one occupied port, regardless of
advice, and extensions preserve ports. Each failed merge consumes two
occupied ports that cannot be used again. Hence there are at most
$\lfloor a/2\rfloor$ failures and
\[
 |\Trust(\text{cycle})|\ge m-\lfloor a/2\rfloor
 \ge m-a/2\ge3m/4,
\]
using $m=2a+b$. Since $\OPT=m$ for each cycle, summing proves the claim.
The bound permits arbitrary correlations between root choices induced by
the advice table.
\end{proof}

\section{Interpolation and matching upper bounds}
The two algorithms give the endpoints in Table~\ref{tab:endpoints}.
The advantage of cycles persists when advice is arbitrarily wrong, rather
than only when the table is close to a maximum matching.

\begin{table}[ht]
\centering\normalsize
\caption{Endpoints used to attain the exact frontiers. Pairs are
(consistency, robustness); all bounds are multiplicative.}
\label{tab:endpoints}
\begin{tabular}{@{}lccc@{}}
\toprule
Realized class & Advice-free $\Fair$ & Advice rule $\Trust$ & Upper bounds\\
\midrule
Maximum degree two & $(3/4,3/4)$ & $(1,1/2)$ & $r\le3/4$, $c+r\le3/2$\\
2-regular & $(7/8,7/8)$ & $(1,3/4)$ & $r\le7/8$, $c+r\le7/4$\\
\bottomrule
\end{tabular}
\end{table}

\paragraph{Achievability.}
Write $\beta=3/4$ on $\D$ and $\beta=7/8$ on $\C$.
Corollary~\ref{cor:fair} gives $(\beta,\beta)$, while
Lemmas~\ref{lem:advice} and~\ref{lem:robust} give $(1,2\beta-1)$.
For $c\le\beta$, use $\Fair$. For $c\ge\beta$, choose $\Trust$ once
before processing the input with probability
\[
 p=\frac{c-\beta}{1-\beta},
\]
and use $\Fair$ otherwise. For every fixed input and table, the mixture
has consistency at least $(1-p)\beta+p=c$ and robustness at least
$(1-p)\beta+p(2\beta-1)=2\beta-c$.

\paragraph{Path upper bounds.}
For completeness, specialize the known Z-path obstruction
of~\cite{CGLB24,JM24} to the matching-table domain. Take $U=\{0,1\}$,
online order $x,y$, and $N(x)=\{0,1\}$ in two worlds. Complete with
$N_0(y)=\{0\}$ and $N_1(y)=\{1\}$. Both graphs belong to $\D$ and have
optimum two. With a common advice table, the first action is the same
when the algorithm's random tape is coupled across worlds. If it matches
$x$, the sum of the best final cardinalities is three; if it abstains,
the sum is two. Thus the sum of expected outputs is at most three.

Use empty advice in both worlds. Robustness implies $4r\le3$.
Alternatively use the same table $A(x)=0,A(y)=1$ in both. This table is
optimal in world 1 and allowed in world 0, so $2c+2r\le3$.

\paragraph{Cycle upper bounds with a common table.}
Use the eight-cycle pair of F\`abrega and Mu\~noz~\cite{FM18},
Proposition~1. Take $U=\{0,1,2,3\}$ and online order $x_1,x_2,y_1,y_2$.
The first two neighborhoods are $N(x_1)=\{0,1\}$ and
$N(x_2)=\{2,3\}$ in both worlds. The remaining neighborhoods are
\[
\begin{array}{c|cc}
 &N(y_1)&N(y_2)\\ \hline
G_0&\{0,2\}&\{1,3\}\\
G_1&\{0,3\}&\{1,2\}.
\end{array}
\]
Each graph is a connected simple 2-regular bipartite graph with optimum
four. Fix a common advice table and couple the random tape. If the two
early vertices both match, one free offline vertex remains in each early
pair. In exactly one world those two free vertices belong to the same late
neighborhood, allowing at most one more match; in the other world at most
two more matches are possible. The sum of final cardinalities is at most
seven. With one or zero early matches the sums are at most six or four.
This argument allows abstention, correlated choices, and arbitrary later
strategies. Consequently the sum of expected outputs is at most seven.

Empty common advice gives $8r\le7$. For the consistency bound, use
\[
 (A(x_1),A(x_2),A(y_1),A(y_2))=(0,2,3,1).
\]
It is a perfect matching of $G_1$ and an allowed injective table for $G_0$.
Only the pair $(y_1,3)$ is absent in $G_0$; $(y_2,1)$ is present in both
worlds. Applying consistency in $G_1$ and robustness in $G_0$ gives
$4c+4r\le7$. These fixed input pairs prove both upper bounds in
Theorem~\ref{thm:frontiers}; the preceding mixture attains them.

\section{Checks and further questions}
The accompanying standard-library checker directly enumerates root coin
strings, maximum matchings, and injective advice tables. It verifies both
component expectation formulas on paths and cycles through six online
vertices, and checks all labelled equal-part graphs of size three under
both degree bounds and all online orders. On these inputs it tests every
maximum matching as advice; on regular inputs it also tests every allowed
table, including absent-edge advice. It checks the upper gadgets by exact
optimal continuation. These finite checks supplement the analytic proofs.

The formulas resolve the matching-table tradeoff under the stated oblivious
input model. Further questions include adaptive adversaries, weighted
matching, and a larger offline degree bound. The boundary-bit invariant and
the decreasing-arrival proof both use the second appearance as the last
possible appearance, so extensions to those models need additional ideas.

\paragraph{AI assistance.}
Codex assisted with proof development, exact computations, literature
comparison, and drafting. A separate model-context pass checked the proofs
and finite certificates.
